\documentclass[11pt]{article}
\usepackage[T1]{fontenc}
\usepackage{lmodern,amsmath,amssymb,amsthm}
\usepackage[margin=1in]{geometry}
\usepackage[hidelinks]{hyperref}
\newtheorem{theorem}{Theorem}
\newtheorem{corollary}[theorem]{Corollary}
\title{Real origin fibres of a modular parametrization and the central analytic rank}
\author{\small \href{mailto:admin@cosmosgeek.org}{admin@cosmosgeek.org}\qquad \href{mailto:b321014@iiit-bh.ac.in}{b321014@iiit-bh.ac.in}}
\date{Working paper for specialist review}
\begin{document}\raggedbottom\emergencystretch=2em\maketitle
\begin{abstract}
For a modular parametrization of an elliptic curve over $\mathbb Q$,
an analytic rank $a\ge2$ forces at least $a-1$ distinct crossings of zero
by its real elliptic logarithm between the cusps $0$ and $\infty$.
Consequently each sufficiently close real point has at least $a+1$ real
preimages. We obtain one-sided central-rank tests by counting real roots
of a polynomial on the origin fibre. For positive root number an empty
open half-arc above the origin, together with $L(E,1)=0$, certifies
analytic rank two. The argument is an application of Mellin moments
and sign changes, classical for the differential itself.
\end{abstract}
Let $E/\mathbb Q$ have conductor $N$, normalized weight-two newform
$f$, and root number $w$. Fix $\phi:X_0(N)\to E$ over $\mathbb Q$ with
$\phi(\infty)=O$ and $\phi^*\omega_E=c\,2\pi i f(z)\,dz$ for a real
$c\ne0$. Set $a=\operatorname{ord}_{s=1}L(E,s)$. We write
\[
 y(t)=e^t/\sqrt N,\qquad X(t)=2\pi\int_{y(t)}^\infty f(iy)\,dy .
\]
In the real uniformization, $cX(t)$ is a lift of $\phi(iy(t))$.
The normalizations of $X$ and the sign below are fixed by $a_1(f)=1$.

\begin{theorem}\label{th:fibre}
If $a\ge2$, at least $a-1$ distinct noncuspidal real algebraic points
on the imaginary geodesic map to $O$. At each of these points the
parametrization crosses $O$ on its real locus. Every point of
$E(\mathbb R)$ sufficiently close to $O$ has at least $a+1$ distinct
real preimages. In particular $\deg\phi\ge a+1$.
\end{theorem}
\begin{proof}
Assume $L(E,1)=0$. Then $X(t)$ tends to zero at both ends of the
real axis. At $+\infty$ this follows from the Fourier expansion,
and at $-\infty$ it follows from the central value and the Fricke
involution. At both ends it decays faster than $e^{-A|t|}$ for every
fixed $A>0$. The Mellin identity and integration by parts give the
entire identity
\[
 D(u):=N^{u/2}(2\pi)^{-u}\Gamma(1+u)L(E,1+u)
   =\int_{\mathbb R}e^{ut}(-X'(t))\,dt
   =u\int_{\mathbb R}e^{ut}X(t)\,dt .
\]
The displayed product is interpreted by continuation. Its prefactor
is $1$ at zero. Thus
$\int_{\mathbb R} t^jX(t)\,dt=0$ for $0\le j\le a-2$.
The function $X$ is real analytic and positive for all sufficiently
large $t$. Fricke symmetry gives $X(-t)=-wX(t)$, so its sign is also
fixed in a left tail. If it had just $h<a-1$ sign changes at $t_i$,
the continuous function $X(t)\prod_{i=1}^h(t-t_i)$ would have constant
sign and be nonzero somewhere. Its strictly nonzero integral
contradicts the moment conditions. Each crossing of zero of $X$ has
odd multiplicity and maps to $O$.

The open imaginary axis embeds in $X_0(N)$. In fact a matrix
$\bigl(\begin{smallmatrix}A&B\\ C&D\end{smallmatrix}\bigr)$ in
$\mathrm{SL}_2(\mathbb Z)$ taking $iy$ to a point on this axis
satisfies $AC y^2+BD=0$. If all four entries are nonzero, $AD$ and
$BC$ must have opposite signs, contradicting $AD-BC=1$. The only
remaining possibilities are $\pm I$ or an antidiagonal matrix with
$|C|=1$, the latter excluded by $N>1$. The fibre above a rational
point is finite over $\mathbb Q$, so these real points are algebraic.

Both cusps map to $O$ and are distinct. Both are unramified:
the first coefficient in the expansion at $\infty$ is one and
Fricke interchanges the cusps. Each of the $a-1$ odd crossings maps
its real neighborhood onto a real neighborhood of $O$, as does each
unramified cusp on the real modular curve. The neighborhoods in the
source can be chosen pairwise disjoint. This proves the assertion.
\end{proof}

\begin{corollary}\label{cor:half}
Suppose $L(E,1)=0$. If $k$ is the number of points of
$\phi^{-1}(O)$ on $iy$, $y>1/\sqrt N$, then
\[
 a\le 2+2k\quad(w=+1),\qquad a\le 1+2k\quad(w=-1).
\]
In the first case, if $k=0$ then $L''(E,1)>0$. In the second case,
if $k=0$ then $L'(E,1)>0$.
\end{corollary}
\begin{proof}
For $w=+1$ write $a=2m$. The odd function $X$ has zero moments on
the positive half-axis against $t,t^3,\ldots,t^{2m-3}$. Fewer than
$m-1$ positive crossings would give a contradiction by using
$t\prod_i(t^2-t_i^2)$ as a test polynomial. For $w=-1$ write $a=2m+1$:
use the evenness of $X$ and $\prod_i(t^2-t_i^2)$ for the moments
against $1,t^2,\ldots,t^{2m-2}$. With $k=0$, $X$ is positive
on the entire open positive half-axis. The Mellin identity gives
$L''(E,1)=4\int_0^\infty tX(t)\,dt$ in the even case and
$L'(E,1)=2\int_0^\infty X(t)\,dt$ in the odd case. The Fricke
midpoint is excluded throughout.
\end{proof}

\section*{An algebraic upper certificate}
Use the rational modular unit
\[
 U_N(z)=N^6\frac{\Delta(Nz)}{\Delta(z)},\qquad
 B_\phi(T)=\prod_{\substack{x\in\phi^{-1}(O)_{\mathrm{red}}\\
                             x\ {\rm noncuspidal}}}(T-U_N(x)).
\]
The product is over geometric points, so its coefficients are rational
by Galois invariance. The unit has no poles at the indicated points.
In the positive real $q=e^{-2\pi y}$ coordinate it equals
\[
 N^6q^{N-1}\prod_{n\ge1}(1+q^n+\cdots+q^{n(N-1)})^{24}.
\]
It is strictly decreasing as $y$ increases. The transformation law of
$\Delta$ gives $U_N(-1/(Nz))=U_N(z)^{-1}$, so $U_N$ maps the
strict upper half of the geodesic bijectively onto $(0,1)$.
The number of distinct real roots of $B_\phi$ in $(0,1)$ is therefore
an upper bound for $k$ in Corollary~\ref{cor:half}. Extraneous roots
can come from other branches. The empty interval criterion involves
only a finite fibre and real root isolation once an explicit algebraic
model of $\phi$ has been obtained, for example by the methods
of Griffin--Hales~\cite{GH}.

\section*{Real versus rational branches}
Suppose in addition that $\phi$ comes from the optimal elliptic quotient
of $J_0(N)$ and $E(\mathbb Q)$ contains a point $P$ of infinite order.
Then there are rational points arbitrarily close to $O$ for which
all the preimages constitute a single Galois orbit. Indeed $\phi$
does not factor geometrically through a nontrivial \'etale isogeny
of $E$. Such a factorization would factor the optimal Jacobian
quotient through that isogeny, contradicting connectedness of its
kernel. Equivalently its pullbacks by multiplication maps are
geometrically irreducible (an unconnected pullback gives a nontrivial
intermediate \'etale quotient). The Hilbert irreducibility theorem
for covers of abelian varieties~\cite[Theorem 1.4]{CDJLZ} gives a
finite-index coset in $\langle 2P\rangle$ on which the fibres are
integral. Each such coset is dense in the real identity circle.
Intersect it with the neighborhood in Theorem~\ref{th:fibre}, omitting
the finitely many branch values. Many real preimages can therefore
belong to one rational orbit even on a single cyclic subgroup.
The sign-change method of Mazur and Swinnerton-Dyer, as presented by
Delaunay~\cite{D}, counts changes of sign of the newform along the
modular path. The fibre count here applies the argument to its
primitive and counts crossings of a specified elliptic-curve value.
\section*{A degree refinement for even central order}
If $w=+1$ and $a\ge2$, the real lift $X$ is odd, so its zero at
$t=0$ has odd order $e$, the ramification index of $\phi$ at
$z_*=i/\sqrt N$. The lattice $\mathbb Z+\mathbb Z z_*$ has CM by
$\mathbb Z[\sqrt{-N}]$, of discriminant $-4N$: the primitive quadratic
equation is $Nz_*^2+1=0$. Its $j$-invariant therefore has degree the
order class number $h(-4N)$ by the main theorem of complex multiplication.
Since the $j$-map is rational, the orbit of $z_*$ on $X_0(N)$ contains
at least this many points, each in the origin fibre with multiplicity
$e$, and each fixed by the rational Fricke involution. On the embedded
imaginary path $z_*$ is the only Fricke fixed point. The two cusps
and at least $a-2$ other sign-changing points found above are disjoint
from this orbit. Hence $\deg\phi\ge e\,h(-4N)+a$. For the CM description
of Fricke fixed points see \cite[\S5]{Oz}.

\section*{Twisting the real path by a unitary character}
This gives a control on a possible use of the real branches.
Write $Y(y)=cX(\log(y\sqrt N))$. Then $Y'=-2\pi c f(iy)$,
$Y(\infty)=0$ and $Y(0)=cL(E,1)$. On $\mathbb C/\Lambda_E$, a
logarithm $u$ determines the flat unitary character of a period
$\lambda$ given by $\exp(2\pi i\,\operatorname{Im}(\overline u\lambda)/
\operatorname{area}(\mathbb C/\Lambda_E))$. Along the real path its
transport has the form $\exp(aY(y))$. If the marked point is twice
a rational point, it lies in the identity real component and admits
a real logarithm, giving constant transport. For arbitrary $a$ the
central integral is still determined by the endpoints:
\[
 \int_0^\infty f(iy)e^{aY(y)}\,dy
   =\frac{e^{a cL(E,1)}-1}{2\pi c a},
\]
with the continuous value at $a=0$. This follows by one integration.
When $L(E,1)=0$, set
$J(s,a)=\int_0^\infty f(iy)y^{s-1}e^{aY(y)}\,dy$.
A second integration by parts gives
\[
 \partial_a J(s,0)=
 \frac{s-1}{4\pi c}\int_0^\infty Y(y)^2 y^{s-2}\,dy .
\]
The last integral is finite and strictly positive at $s=1$. Thus this
first character coefficient has a simple central zero for every such
curve. In particular its order is unchanged when the original newform
has a higher-order central zero. For modular-symbol character
deformation see \cite{PR}.

\section*{Two exact readings of the analytic order}
The fibre criterion above bounds the central order using a finite
algebraic fibre. For comparison, consider two identities whose inputs
are themselves $L$-functions. First take $w=+1$, $L(E,1)=0$, and let
$\pi$ be the corresponding representation of $\mathrm{PGL}_2$ with
unitary central parameter $1/2$. Schmidt's Saito--Kurokawa packet
with empty exceptional set contains a cuspidal member $\Pi$ under
precisely this vanishing hypothesis \cite[Thm.~3.1]{SK}. Its spin
parameter is the sum of that of $\pi$ and $|\cdot|^{\pm1/2}$
\cite[Prop.~5.5]{Gan}. With the unitary normalization
$\Lambda_\pi(t)=\Lambda_E(t+1/2)$, where
$\Lambda_E(s)=(\sqrt N/(2\pi))^s\Gamma(s)L(E,s)$, it follows that
\[
 \Lambda_{\rm spin}(\tfrac12+u,\Pi)
 =\Lambda_E(1+u)Z(u)Z(1+u),\qquad
 Z(t)=\pi^{-t/2}\Gamma(t/2)\zeta(t).
\]
The residues of $Z$ at $0$ and $1$ are $-1$ and $+1$. Thus
\[
 \Lambda_{\rm spin}(\tfrac12,\Pi)=-\tfrac12\Lambda_E''(1)
   =-\frac{\sqrt N}{4\pi}L''(E,1).
\]
This is a specialization of the existing packet formula. The bare
finite product has only one pole available to remove central order;
the second comes from completion. The cuspidal member here need not be
holomorphic. Its nonzero existence does not require $L''(E,1)\ne0$.
In particular, existence of the representation supplies no test for
the vanishing of this spin value.

For the second identity, take an $S_3$ extension $D/\mathbb Q$, a
cubic subfield $M$, and its quadratic subfield $K$. Put
$F_J(t)=L(E/J,1+t)$, $a_J=\operatorname{ord}_{t=0}F_J(t)$ and use
meromorphic germs. Artin formalism gives
$F_D F_{\rm Q}^2=F_M^2F_K$: the permutation characters for
$D,M,K$ are respectively the regular representation,
$\mathbf1+\mathrm{std}$, and $\mathbf1+\mathrm{sgn}$. The relation
implies both equality of orders and equality of initial coefficients.
Consequently
\[
 \lim_{t\to0}\frac{F_D(t)F_{\rm Q}(6t)^2}
 {F_M(2t)^2 F_K(3t)}
 = 2^{2a_{\rm Q}-2a_M}3^{2a_{\rm Q}-a_K}.
\]
If $L(E\otimes\chi_K,1)\ne0$ the exponent of three is $a_{\rm Q}$.
Choose such a nontrivial twist by the nonvanishing theorem of
Friedberg--Hoffstein \cite{FH}. For completeness any corresponding
$K=\mathbb Q(\sqrt d)$ occurs in an $S_3$ extension: put
$U=1+3dT^2$ and $g(X)=X^3-3UX-2U$. The generic cubic is irreducible
over $\mathbb Q(T)$, since the defining relation solves for
$T^2=(X-2)(X+1)^2/(3d(3X+2))$, which has an odd-order zero and pole
in $\mathbb Q(X)$; Gauss's lemma then applies in the two variables.
Its discriminant is $324d\,U^2T^2$. Hilbert irreducibility gives an
irreducible specialization with nonzero discriminant, whose splitting
field therefore has quadratic subfield $K$.

There is a restriction on interpreting the rescaled quotient as a
new ordinary $L$-function. For finite $L$-functions with finitely many
Euler factors omitted, take the quotient in the last display and remove
its $F_D(t)$ factor. In the variable $s=1+t$ all remaining Euler
substitutions are $s\mapsto1+k(s-1)$ with $k=2,3,6$. Thus every
good linear Euler coefficient vanishes. At a prime $q$ splitting
completely in $D$, the coefficient of $T^2$ with $T=q^{-s}$ equals
$-6q a_q(E)$: the only contribution at that degree is from the inverse
square of $F_M(2t)$, with three equal local factors. Infinitely many
such $q$ have $a_q(E)\ne0$, by simultaneously splitting in the odd
division field $\mathbb Q(E[\ell])$ for any fixed $\ell\ge3$; then
$a_q(E)\equiv2\bmod\ell$. If an ordinary virtual finite-dimensional
continuous $\ell$-adic representation had these local factors, its
trace would vanish at almost all Frobenius classes, so its
semisimplified virtual character would be zero by density and
Brauer--Nesbitt. Its good local factors would then all be one, contrary
to the quadratic coefficient. This obstruction concerns the proposed
Euler realization and leaves the germ identity intact. For the
ordinary Artin relations on elliptic $L$-functions and regulators see
\cite{DD}.
\subsection*{A circular mean for a shifted CM lattice}
Suppose $E/\mathbb Q$ has complex multiplication by the full ring
$\mathcal O$ of an imaginary quadratic field of class number one.
Write $w=|\mathcal O^\times|$. Choose the Hecke character of $E$ so that
$\psi((\alpha))=\epsilon(\alpha)\alpha$ and $L(\psi,s)=L(E,s)$.
Extend its finite character $\epsilon$ by zero on nonunits modulo its
conductor. Then $\epsilon(v)=v^{-1}$ for $v\in\mathcal O^\times$, so
its mean over a full residue system is zero. For $z$ near zero set
\[
 H(u,z)=\frac1w\sum_{\epsilon(\alpha)\ne0}
       \epsilon(\alpha)\frac{\alpha+z}{|\alpha+z|^{2+2u}}.
\]
The sum starts in $\Re u>1/2$, where $H(u,0)=L(E,1+u)$ by summing
over generators of principal ideals. For the CM lattice-series
framework see F. Rodriguez Villegas and D. Zagier,
\href{https://people.mpim-bonn.mpg.de/zagier/files/mpim/92-68/fulltext.pdf}
{\emph{Square roots of central values of Hecke L-series}}.

Insert the Mellin integral for $|\alpha+z|^{-2-2u}$.
Poisson summation on the conductor-lattice cosets gives exponential
decay as the heat parameter tends to zero: the zero dual mode
vanishes by the mean of $\epsilon$. The same holds with the linear
numerator removed or after any fixed number of spatial derivatives.
Decay at the other endpoint is exponential locally uniformly for
$|z|$ smaller than the distance to the nearest contributing lattice
point. Division by $\Gamma(1+u)$ thus gives continuation near $u=0$
with bounded spatial derivatives on a fixed smaller disk.

For $\Delta=4\partial_z\partial_{\bar z}$, differentiation in the
absolute-convergence region and continuation give
\[
 \Delta H(u,z)=4u(1+u)H(u+1,z).
\]
Likewise $\partial_z H$ is $-u$ times the continued scalar lattice
sum with denominator $|\alpha+z|^{2+2u}$, including the same
$1/w$ normalization. Consequently $H(0,z)$ is antiholomorphic.
It is invariant under $z\mapsto vz$ for every unit $v$.
If $L(E,1)=0$, the origin in the spatial variable has multiplicity
exactly $w$: the coefficient of $\bar z^w$ equals
\[
 \frac{(-1)^w}{w}\sum_{\epsilon(\alpha)\ne0}
       \frac{\epsilon(\alpha)}{\bar\alpha^{\,w+1}} .
\]
This differentiated sum is absolutely convergent. It is $(-1)^w$
times the Hecke series at $s=w+1$ of the character on ideals
$(\alpha)\mapsto\epsilon(\alpha)\alpha^{w+1}$. Its Euler product is
absolutely convergent and nonzero there, since $(w+1)/2>1$.

If $L(E,1)=L'(E,1)=0$ and $\lambda\ge0$ is fixed, then
\[
 \lim_{u\downarrow0}\frac1{2\pi u^2}\int_0^{2\pi}
       H\!\left(u,\sqrt{\lambda u}\,e^{i\theta}\right)\,d\theta
 =\tfrac12 L''(E,1)+\lambda L(E,2).
\]
Indeed the circular Taylor mean at radius $\rho$ is
$H+(\rho^2/4)\Delta H+(\rho^4/64)\Delta^2 H+O(\rho^6)$, with all
three terms evaluated at $(u,0)$. Here $\Delta^2H(u,0)=O(u)$ by
two applications of the differential identity. The displacement
term survives for every positive $\lambda$ even if $L''(E,1)=0$.


\begin{thebibliography}{9}\raggedright
\bibitem{D} C. Delaunay, \emph{Critical and ramification points of the modular
parametrization of an elliptic curve}, J. Th\'eor. Nombres Bordeaux
\textbf{17} (2005), 109--124. \href{https://doi.org/10.5802/jtnb.480}{Paper}.
\bibitem{GH} M. Griffin and J. Hales, \emph{Divisors of modular parametrizations
of elliptic curves}, \href{https://arxiv.org/abs/2003.01675}{arXiv:2003.01675}.
\bibitem{CDJLZ} P. Corvaja, J. Demeio, A. Javanpeykar, D. Lombardo and
U. Zannier, \emph{On the distribution of rational points on ramified covers
of abelian varieties}, \href{https://arxiv.org/abs/2011.12840}{arXiv:2011.12840}.
\bibitem{Oz} E. \"Ozman, \emph{Points on quadratic twists of $X_0(N)$},
\href{https://arxiv.org/abs/0911.4537}{arXiv:0911.4537}.
\bibitem{PR} Y. Petridis and M. Risager, \emph{Arithmetic statistics of
modular symbols}, \href{https://arxiv.org/abs/1703.09526}{arXiv:1703.09526}.
\bibitem{SK} R. Schmidt, \emph{The Saito--Kurokawa lifting and functoriality},
\href{https://www.math.uni-sb.de/ag/schulze/rschmidt/papers/skfunc.pdf}{paper}.
\bibitem{Gan} W. T. Gan, \emph{The Saito--Kurokawa space of PGSp$_4$ and its transfer to inner forms},
\href{https://math.nyu.edu/~tschinke/princeton/papers/.aim05/submitted/gan.pdf}{paper}.
\bibitem{FH} S. Friedberg and J. Hoffstein, \emph{Nonvanishing theorems
for automorphic $L$-functions on $\mathrm{GL}(2)$},
\href{https://annals.math.princeton.edu/1995/142-2/p04}{paper}.
\bibitem{DD} T. Dokchitser and V. Dokchitser,
\emph{On the Birch--Swinnerton-Dyer quotients modulo squares},
\href{https://annals.math.princeton.edu/wp-content/uploads/annals-v172-n1-p11-p.pdf}{paper}.
\end{thebibliography}\end{document}